%!TEX Root = ../main.tex
\chapter{Rings}
\label{chap:rings}
The origins of commutative algebra goes back
to about 150 years ago, when the following subjects
are studied.
\begin{itemize}
    \item \textbf{Algebraic Number Theory}, which
	studies about a specific kind of commutative rings
	generated from \vocab{number fields}.
	(e.g., $\mathbb Q$, $\mathbb Z[i]$, $\mathbb Z_p$,
	$\mathbb Z/p\mathbb Z$, etc.)
    \item 
    \textbf{Algebraic Geometry},
    which studies about polynomial rings
    $k[x]$, $k[x_1,\dots,x_n]$, and their quotients
    (where $k$ is a field).
    The most important casae is when $k=\mathbb C$,
    but modern algebraic geomety have extended to studying
    polynomial equations over arbitrary field too.
\end{itemize}
However, people then realize that these two subjects
have the same foundation.
In particular, theorems about Algebraic Number Theory and
theorems in Algebraic Geometry have a lot in common,
and the common foundation
--- properties of commutative rings --- establishes
\vocab{commutative algebra}.
The central object of commutative algebra is ring,
which we introduce in this chapter.

\section{Rings and Ideals}
In this section, we set up some basic concepts of rings and ideals.
\subsection{Ring Homomorphisms}
The most basic object we are going to study
is (commutative) rings. A commutative \vocab{ring} $A$ has operations addition $+$ and multiplication $\cdot$, which 
must satisfies the following:
\begin{itemize}
    \item addition forms an abelian group,
    i.e., it must be 
    \begin{itemize}
    \item associative: $a+(b+c)=(a+b)+c$ for all $a,b,c\in A$.
    \item commutative: $a+b=b+a$ for all $a,b\in A$.
    \item identity: there exists $0\in A$ such that $a+0=0+a=a$
    for all $a\in A$.
    \item inverse: for all $a\in A$, there exists $(-a)\in A$ 
    such that $a+(-a)=(-a)+a=0$.
    \end{itemize}
    \item multiplication is 
    \begin{itemize}
    \item associative: $(a\cdot b)\cdot c = a\cdot (b\cdot c)$
    for all $a,b,c\in A$.
    \item commutative: $a\cdot b = b\cdot a$ for all $a,b\in A$.
    \item has identity $1\in A$ such that $a\cdot 1 = 1\cdot a
    = a$ for all $a\in A$.
    \end{itemize}
    \item Multiplication distributes over addition:
    $a\cdot (b+c)= (a\cdot b) + (a\cdot c)$
    and $(b+c)\cdot a = (b\cdot a) + (c\cdot a)$
    for all $a,b,c\in A$.
\end{itemize}
We allow $0=1$, in which case $A=\{0\}$
is a \vocab{zero ring}.

The word ``commutative'' in commutative rings 
indicates that multiplication is commutative.
Noncommutative rings are not the focus of these notes.
\begin{center}
    \fbox{In these notes, we assume all rings are commutative and contain $1$.}
\end{center}
\begin{example}[Basic rings]
The most basic example of ring is integers $\ZZ$. 
Furthermore, $\QQ$, $\RR$, and $\CC$ are all rings.
Integers modulo $n$, $\ZZ/n\ZZ$, form a ring.
\end{example}
\begin{example}[Constructing new rings]
Here are some ways to construct a new ring.
\begin{itemize}
\item Given a ring $A$, one can adjoin a free element $x$
to get a \vocab{polynomial ring} $A[x]$, which is 
described by 
$$A[x] = \{a_0 + a_1x + \dots + a_nx^n : a_0,\dots,a_n\in A\}.$$
\notation{$A[x]$}{polynomial ring}
Similarly, one can adjoin arbitrarily many elements 
to get rings such as $A[x_1,\dots,x_m]$, 
which is a ring of polynomials in $m$ variables
with coefficients in $A$.
\item Given two rings $A$ and $B$, one construct 
the \vocab{product ring}
$$A\times B = \{(r,s) : r\in A, s\in B\},$$
\notation{$A\times B$}{product ring}
where addition and multiplication are done component-wise.
\end{itemize}
\end{example}
A \vocab{ring morphism} between two rings $A$ and $B$
is a map $\phi : A\to B$,
that satisfies $\phi(x+y) = \phi(x)+\phi(y)$, $\phi(xy) = \phi(x)\phi(y)$, and $\phi(1)=1$ for all $x,y\in A$.
In particular,
\begin{center}
    \fbox{We assume that all ring morphisms 
    preserve $1$.}
\end{center}
\begin{example}[Ring morphisms]
    We have the following examples:
    \begin{itemize}
	\item For any ring $A$, there exists a unique ring
    homomorphism $f : \mathbb Z\to A$
    by
    $$f(n) = \underbrace{1+\dots+1}_n.$$
    \item For any ring $A$, there exists a 
	unique ring morphism $f : A\to \{0\}$.
    \item If $p$ is a prime, and suppose
	$\underbrace{1+\dots+1}_p = 0$
	(i.e., $A$ has characteristic $p$),
	then we have the \vocab{Frobenius morphism}
	$F : A\to A$ by
	$F(a) = a^p$.
	(Proof: multiplication is obvious.
	For addition, $(a+b)^p = a^p + \binom p1 a^{p-1}b
	+ \binom p2 a^{p-2}b^2 + \dots + b^p = a^p+b^p$.)
    \end{itemize}
\end{example}
\subsection{Ideals}
Ideals are a subset of ring that has special properties.
Most importantly, it forms a quotient ring, which we 
describe in the next subsection.
A subset $I\subseteq A$ is an \vocab{ideal} if and only if it
\begin{itemize}
    \item is closed under addition:
	if $a,b\in I$, then $a+b\in I$; and
    \item absorbs multiplication:
	if $a\in I$ and $r\in A$, then $ra\in I$.
\end{itemize}
Notice that the second condition is stronger than 
being closed under multiplication.
\begin{example}
    We have the following examples of ideals.
    \begin{itemize}
	\item For any integer $n$, 
        $(n) = n\ZZ\subseteq \ZZ$ is an ideal.
        Moreover, all ideals of $\ZZ$ are of this form.
	\item $\{0\}$ and $A$ itself is an ideal of $A$.
	    If $A$ is a field, these two
	    are the only ideals of $A$.
	\item For each $a\in A$,
	    we have the ideal 
	    generated by $a$ defined by
	    $(a) = \{ar : r\in A\}$.
	    These ideals are called \vocab{principal ideals}.
	    For $\mathbb Z$, every ideal is principal.
	\item More generally, given some subset $T\subseteq A$,
	    we can define the ideal generated by $T$
	    as the smallest ideal containing $T$, which is
	    $$(T) = \left\{\sum_{i=1}^m a_it_i : 
	    a_i\in A, t_i\in T\right\}.$$
            Ideals generated by $\{x_1,\dots,x_n\}$
            is often written as $(x_1,\dots,x_n)$.
            \notation{$(T)$}{ideal generated by set $T$}
            \notation{$(x_1,\dots,x_n)$}{ideal generated by $x_1,\dots,x_n$}
    \end{itemize}
\end{example}
Another way to construct ideals is to take the kernel.
Given a ring homomorphism $f : A\to B$, then
$$\Ker f= f^{-1}(\{0\})\subseteq A$$
\notation{$\Ker f$}{kernel of ring homomorphism}
is always an ideal of $A$.
This is in fact a special case of
a more general construction called contraction.
\begin{definition}[Contraction and Extension]
    Consider a ring homomorphism
    $f : A\to B$. 
    \begin{itemize}
    \item For each ideal $J\subseteq B$,
    the \vocab{contraction} 
    $f^{-1}(J)\subseteq A$ is automatically an ideal.
    \item For each ideal $I\subseteq A$,
    the \vocab{extension} $IB$
    \notation{$IB$}{extended ideal}
    is the ideal \textbf{generated by} the image $f(I)$.
    \end{itemize}
\end{definition}
Finally, we have the following operations of ideals.
\begin{definition}[Ideal operations]
    Let $I,J$ be an ideal of a ring $A$. Then,
    \begin{itemize}
    \item $I\cap J$ is an ideal.
    \item $I+J = \{i+j : i\in I, j\in J\}$ is an ideal.
    \notation{$I+J$}{ideal addition}
    \item $IJ$ is the ideal generated by $\{ij : i\in I, j\in J\}$.
    \notation{$IJ$}{ideal multplication}
    \end{itemize}
    They are analogue of $\lcm$, $\gcd$, and product, respectively.
    For example, in $\ZZ$, we have 
    $(4)\cap (6) = (12)$, $(4)+(6)=(2)$,
    and $(4)(6) = (24)$.
\end{definition}
\begin{remark}[Origin of ideals]
The first usage of the word \textbf{ideal}
is by Dedekind, who is a algebraic number theorist.
He coined the word \textbf{ideal numbers}
to describe the phenomenon of
unique factorization in rings of algebraic integers.

To elaborate, people at that time are interested
in solving the equation $x^n+y^n=z^n$ (Fermat's last theorem)
where $x,y,z\in \ZZ$.
One idea (led by Kummer) to solve
this problem is to adjoin $\zeta = e^{2\pi i/n}$, then
one can rewrite this equation to 
$$(x+y)(x+\zeta y)\dots (x+\zeta^{n-1}y)=z^n.$$
Then, people work with the ring $\mathbb Z[\zeta]$
and, assuming that it's UFD (i.e., has unique factorization property), able to deduce that each factor $x+y$, $x+\zeta y$,
\dots, $x+\zeta^{n-1}y$ must be an $n$-th power
(or close to), and then make a contradiction.

However, the ring of algebraic integers is not necessarily a UFD. For example, in the ring $\ZZ[\sqrt{-5}]$, we have
$2\cdot 3 = (1+\sqrt{-5})(1-\sqrt{-5})$,
and neither $2$, $3$, $1+\sqrt{-5}$, nor $1-\sqrt{-5}$
can be factored further in $\ZZ[\sqrt{-5}]$.
Thus, $\ZZ[\sqrt{-5}]$ is not a UFD.
In the case of Fermat's last theorem, the smallest $n$
for which unique factorization property fail is $n=23$.

In the world of ideals, both $2$ and $3$ breaks up in $\ZZ[\sqrt{-5}]$.
Indeed, in ideals,
$$(2) = (2, 1+\sqrt{-5})^2,\quad (3) = (3, 1+\sqrt{-5})
(3, 1-\sqrt{-5}),$$
and unique factorization is restored, except that
the ideal $(2, 1+\sqrt{-5})$ is no longer principal.
In \Cref{sec:dedekind}, we will discuss a broad class
of ring that unique factorization of ideals hold.
\end{remark}
\subsection{Quotient Rings}
We now describe the most important way to generate new rings,
called \vocab{quotient ring}.
Given a ring $A$ and an ideal $I\subseteq A$,
then we define the quotient ring $A/I$
\notation{$A/I$}{quotient ring}
as the ring of cosets 
$$\{[a] := a+I : a\in A\}.$$
We define the operations by
$[a+b] = [a]+[b]$ and $[ab] = [a][b]$ 
as expected. 
This map is well-defined.
Indeed, if $a-a'\in I$, then e.g., $ab-a'b\in I$,
so $[ab] = [a'b]$.

We then also define the surjective ring homomorphism
$\pi : A\to A/I$ by $\pi(a) = [a]$.
The kernel of this map is $I$.
This map satisfies the following universal property:
suppose $B$ is another ring and a homomorphism
$f : A \to B$ such that $f(I) = 0$, then
there exists a unique ring homomorphism
$\overline f : A/I\to B$ such that
the following diagram commutes:
\begin{center}
\begin{tikzcd}
    A \arrow[rr, "\pi"]
    \arrow[rd, "f"]
    &
    & A/I
    \arrow[ld, "\exists ! \overline f"] \\
    & B
\end{tikzcd}
\end{center}
(Proof: we are forced to have $\overline f([a]) = f(a)$
, so $f$ is unique. Moreover, this is well-defined
because $f$ sends $I$ to $0$.)

By the usual argument (see \Cref{sec:universal}),
this universal property classify $\pi$ and $A/I$
up to isomorphism.

We now state an important theorem characterizing 
ideals of $A/I$.
\begin{theorem}[Correspondence Theorem]
    \index{correspondence theorem}
    \label{thm:correspondence}
    There is a bijection between
    $$\{\text{ideals of } A/I\}
    \ \longleftrightarrow\ 
    \{\text{ideals of } A \text{ containing } I\}.$$
\end{theorem}
For example, the ideals of $\mathbb Z/6\mathbb Z$
are $0$, $\mathbb Z/6\mathbb Z$, $2\mathbb Z/6\mathbb Z$,
and $3\mathbb Z/6\mathbb Z$,
while the ideals of $\mathbb Z$ that contains
$6\mathbb Z$ are $6\mathbb Z$, $\mathbb Z$,
$2\mathbb Z$, and $3\mathbb Z$.
\begin{proof}
    Let $\pi : A\to A/I$ be the quotient projection map.
    The bijection is given by 
    \begin{align*}
    \{\text{ideals of } A/I\}
    \ &\leftrightarrow\ 
    \{\text{ideals of } A \text{ containing } I\} \\
    J' \ &\mapsto\ \pi^{-1}(J') \\
    \pi(J)\ &\mapsfrom\ J 
    \end{align*}
    and it is straightforward to verify that 
    they are inverses.
\end{proof}
\begin{theorem}[First Isomorphism Theorem]
\label{thm:first_isomorphism}
\index{first isomorphism theorem}
Let $\phi : A\to B$ be a surjective ring morphism.
Then, $B\simeq A/\Ker\phi$.
\end{theorem}
\begin{exercise}
Prove the above theorem (\Cref{thm:first_isomorphism}).
\end{exercise}
\section{Properties of Rings}
\subsection{Special Elements}
We have a bunch of definitions of
special elements of a ring.
\begin{definition}[Units, zero divisors, nilpotents]
    Given a ring $A$, an element $a\in A$ is
    \begin{itemize}
    \item a \vocab{unit} if there exists $b\in A$ such that $ab=1$;
    \item a \vocab{zero divisor} if
	there exists $b\neq 0\in A$ such that $ab=0$; and
    \item \vocab{nilpotent} if $a^n=0$ for some $n\geq 1$.
    \end{itemize}
\end{definition}
The set of nilpotent elements turns out to be an ideal.
\begin{proposition}
    \notation{$\Nil A$}{nilradical}
    The set $\operatorname{Nil}A\subseteq A$
    is an ideal of $A$, called the \vocab{nilradical}
    of $A$.
\end{proposition}
\begin{proof}
    Absorbing multiplication is obvious.
    To show that $\Nil(A)$ is closed under addition,
    if $a^m=0$ and $b^n=0$, then the binomial expansion of $(a+b)^{m+n}$ 
    consists of monomials of the form $a^ib^j$.
    At least one of $i\geq m$ and $j\geq n$ must hold,
    so the monomial is divisible by either $a^m$ or $b^n$,
    both of which are $0$, so all monomials vanish,
    which means $(a+b)^{m+n}=0$, so $a+b\in \Nil(A)$.
\end{proof}
Each kind of special element leads to a type of ring.
\begin{definition}[Field, domain, reduced rings]
    A ring $A$ is
    \begin{itemize}
    \item a \vocab{field} if $A\neq 0$ and every $a\neq 0$ is a unit;
    \item a \vocab{domain} if $A\neq 0$ and $0$ is the only zero divisor; and
    \item \vocab{reduced} if $0$ is the only nilpotent element.
    \end{itemize}
\end{definition}
\begin{remark}
    We use the word domain 
    instead of integral domain to avoid confusion caused
    by the word integral as in \Cref{chap:integral}.
\end{remark}
\begin{example}
    \begin{itemize}
    \item $\QQ$, $\RR$, $\CC$, and $\FF_p = \ZZ/p\ZZ$ are fields.
    \item $\mathbb Z$ and $\mathbb C[x]$ are domains.
    (In fact, if $A$ is a domain, then so is $A[x]$,
    so $\CC[x_1,\dots,x_n]$ is a domain.)
    \item $\mathbb C[x,y] / (xy)$ is not a domain
    because in this ring, $xy=0$ but $x,y\neq 0$.
    \item $\CC[x,y]/(xy)$ is reduced.
    \item $\CC[x]/(x^2)$ is not reduced because $x$ is nilpotent.
    \end{itemize}
\end{example}
\begin{exercise}
    What are the zero divisors and nilpotent elements
    of $\mathbb Z/12\mathbb Z$?
\end{exercise}
\subsection{PID and UFD}
\begin{definition}[PID and UFD]
    A domain $A$ is a 
    \begin{itemize}
    \item \vocab{principal ideal domain (PID)}
    if every ideal principal.
    \item \vocab{unique factorization domain (UFD)}
    if every non-zero element $x\in A$ 
    can be written uniquely as a product of irreducible
    elements 
    $$x = p_1\cdots p_n,$$
    where uniqueness is up to multiplication by units,
    i.e., if $x=q_1\cdots q_m$ is another factorization,
    then $m=n$ and there exists a permutation $\pi$
    such that $q_{\pi(i)}$ is equal to $p_i$ times a unit
    for all $i$.

    An element $p\in A$ is \vocab{irreducible} if it is 
    not a unit and not 
    a product of two non-units in $A$.
    \end{itemize}
\end{definition}
We note down two useful properties of PID and UFD proved in
undergraduate algebra. The proofs of these 
are not central to these notes, so we will only sketch it.
\begin{theorem}
    \label{thm:ufd_properties}
    \begin{alphenum}
    \item Any PID is a UFD.
    \item If $A$ is a UFD, then so is $A[x]$
    (hence, by induction, $A[x_1,\dots,x_n]$ is a UFD).
    \end{alphenum}
\end{theorem}
\begin{proof}[Proof Sketch]
\begin{alphenum}
\item First, we need to show that every element $x$
is a product of irreducible elements.
Obviously, we keep factoring out elements.
If this does not terminate, then there would be an infinite
chain $x=a_0, a_1, a_2,\dots$ such that $a_{i+1}$ is a proper divisor of $a_i$, i.e., $(a_0)\subsetneq (a_1)\subsetneq \dots$.
Then, the union $\bigcup_{i\geq 0} (a_i)$ is an ideal,
say $(x)$, then $(x)=(a_i)$, a contradiction.

Next, we show uniqueness. If $p$ is irreducible,
then $(p)$ is maximal, hence prime, so if $p$ divides $ab$,
then $p$ divides either $a$ or $b$.
Thus, if we have two different product of irreducible elements
$x=p_1\cdots p_m = q_1\cdots q_n$, then $p_i$ divides
one of $q_1,\dots,q_n$, hence equal to $q_j$ for some $j$
(up to multiplication by units). Similarly each $q_j$ is equal
to $p_i$ for some $i$. Thus, the multiset of $p_i$ 
and of $q_j$ are equal.
\item 
First, any polynomial in $A[x]$ can be written as 
a product of irreducible polynomials because 
we can repeatedly factor anything that is not 
irreducible, and the degree decreases
until it is constant polynomial, at which point
we can utilize the fact that factoring in $A$ terminates.

To show uniqueness, by the same argument in (a),
it suffices to show that if $f,g,h\in A[x]$
such that $f$ is irreducible and $f\mid gh$,
then $f\mid g$ or $f\mid h$.
To do this, we utilize the following lemma:
``A polynomial $f\in A[x]$ is primitive if and only if
no element in $A$ divides all coefficients in $f$.
If $f,g$ is primitive, then $fg$ is primitive.''
To prove the lemma, suppose that an irreducible factor 
$p$ divides all coefficients of $fg$. Then, $fg=0$ in
$A/(p)[x]$. However, $A/(p)[x]$ is a domain (because $A/(p)$
is a domain), a contradiction.

Let $K$ be the field of fraction of $A$.
If $f,g,h\in A[x]$ such that $f$ is irreducible
(hence primitive) and $f\mid gh$, then 
$f\mid gh$ in $K[x]$, so there exists $d\in A$
such that $\overline f(x) = d\cdot gh/f \in A[x]$.
Scale down so that $\overline f(x)$ is primitive,
then the lemma forces $d=1$, so $f$ divides $gh$ in $A[x]$.
Hence, every irreducible element is prime.
\qedhere
\end{alphenum}
\end{proof}
\begin{example}[PID and UFD]
    \begin{itemize}
    \item $\mathbb Z$, $\mathbb Z[i]$, and $\CC[x]$
    are PID.
    \item $\CC[x,y]$ is a UFD by \Cref{thm:ufd_properties} (b),
    but it is not a PID because ideal $(x,y)$ is not principal;
    if $(x,y)=(f)$, then $x$ and $y$ must be divisible by $f$,
    forcing $f=1$, which is a contradiction.
    \item $\ZZ[\sqrt{-5}]$ is not a UFD
    because $2\cdot 3 = (1+\sqrt{-5})(1-\sqrt{-5})$,
    and $2$, $3$, $1+\sqrt{-5}$, and $1-\sqrt{-5}$
    are all irreducible.
    \end{itemize}
\end{example}

\subsection{Chinese Remainder Theorem}
There is a number-theoretic phenomenon that
for example, $\mathbb Z/(12)\simeq \mathbb Z/(4)
\times \mathbb Z/(3)$.
This fact generalizes to any rings.
\begin{definition}[Coprime ideals]
    Given a ring $A$, two ideals $I,J\subset A$
    are \vocab{coprime} if $I+J=(1)$.
\end{definition}
\begin{theorem}[Chinese Remainder Theorem]
    \label{thm:crt}
    \index{Chinese Remainder Theorem}
    If $I$ and $J$ are coprime, then
    \begin{enumerate}[label=(\alph*)]
    \item $I\cap J = IJ$.
    \item The map $f : A\to A/I\times A/J$ by
	$a\mapsto (a\bmod I, a\bmod J)$ is surjective.
    \item We have
	$A/IJ \simeq A/I\times A/J$.
    \end{enumerate}
\end{theorem}
The proof is the same as number theory class.
\begin{proof}
    The fact that $I$ and $J$ are coprime
    means that there exists $i\in I$ and $j\in J$
    for which $i+j=1$.
    \begin{enumerate}[label=(\alph*)]
    \item The inclusion $IJ\subseteq I\cap J$ is clear.
    Thus, we show that $IJ\supseteq I\cap J$.
	Let $a\in I\cap J$.
	Then, we have
	$$a = a\cdot 1 = a\cdot j + i\cdot a,$$
	and each term is in $I\cdot J$.
    \item Let $[x]\in A/I$ and $[y]\in A/J$.
	Set $a=xj+yi$, and see that
	$a\equiv xj\equiv x\pmod I$ and $a\equiv yi
        \equiv y\pmod J$.
    \item By the first isomorphism theorem
    (\Cref{thm:first_isomorphism}),
	we have $A/\Ker f \simeq A/I\times A/J$.
        By (a),  we have $\Ker f = I\cap J = IJ$,
        so we are done.
	\qedhere
\end{enumerate}
\end{proof}

\section{Primes and Maximal Ideals}
\subsection{Definition and Basic Properties}
\begin{definition}[Prime and Maximal Ideals]
    \begin{itemize}
    \item An ideal $\frakp\subset A$ is a \vocab{prime ideal}
    if $\frakp\neq A$ and $ab\in \frakp$
    implies $a\in \frakp$ or $b\in \frakp$ (or both).
    \item An ideal $\mathfrak m\in A$ is a \vocab{maximal ideal}
    if $\mathfrak m\neq A$ and
    there are no ideals $I$ such that
    $\mathfrak m\subsetneq I\subsetneq A$.
    \end{itemize}
\end{definition}
\begin{proposition}
    \begin{alphenum}
    \item An ideal $\pp\subseteq A$ is prime 
    if and only if $A/\frakp$ is a domain.
    \item An ideal $\mathfrak m\subset A$ is maximal
    if and only if $A/\mathfrak m$ is a field.
    \end{alphenum}
    Consequently, all maximal ideals are prime.
\end{proposition}
\begin{proof}
    \begin{alphenum}
    \item 
    Clear by simply reduce the definition of prime ideal modulo $\pp$.
    \item 
    By correspondence theorem (\Cref{thm:correspondence}),
    $\frakm$ is maximal if and only if $A/\frakm$ has no ideal other than $0$ or $A/\frakm$.
    
    If $A/\mathfrak m$ is a field, then $A/\frakm$
    has only two ideals, because any nonzero ideal $I$
    must contain a nonzero element, which is a unit,
    forcing $I$ to include all elements in $A/\frakm$.
    
    Otherwise, if $A/\mathfrak m$ is not a field,
    then it has a non-unit $a$,
    so $(a)$ is a nontrivial ideal of $A/\mathfrak m$. 
    \qedhere 
    \end{alphenum}
\end{proof}
\begin{example}[Prime and Maximal Ideals]
    We have the following examples.
    \begin{itemize}
    \item $(0)$ is prime if and only if $A$ is a domain.
    \item $(0)$ is maximal if and only if $A$ is a field.
    \item If $A$ is a PID, then all prime ideals of $A$ are 
    $(0)$ and $(f)$ for irreducible element $f$.
    Special cases of this are 
    \begin{itemize}
    \item The prime ideals of $\mathbb Z$
	are $(0)$ and $(p)$, where $p$ is the prime numbers.
	Only $(p)$ are maximal.
    \item The prime ideals of $\CC[x]$ are $(0)$ and 
    $(x-\alpha)$, where $\alpha\in\mathbb C$.
    \end{itemize}
    \item Prime ideals of $\CC[x_1,\dots,x_n]$ are more numerous. For example, 
    $$0\subset (x_1)\subset (x_1,x_2)\subset\dots
	\subset (x_1,\dots,x_n),$$
    are all prime ideals 
    (Proof: $\mathbb C[x_1,\dots,x_n]/(x_1,\dots,x_k)
    = \mathbb C[x_{k+1},\dots,x_n]$, which is a domain).
    On the other hand, only $(x_1,\dots,x_n)$
    is maximal in this chain.

    We will show later in the text (Hilbert Nullstellensatz, 
    \Cref{thm:max_ideal_poly_ring}) that all maximal ideals
    $\mathbb C[x_1,\dots,x_n]$ are in from
    $(x_1-\alpha_1,\dots, x_n-\alpha_n)$,
    where $\alpha_1,\dots,\alpha_n\in\mathbb C$.
    \item Let $k$ be a field and consider the 
    \vocab{formal series ring}
    $$k[[x]] = \left\{\sum_{n=0}^\infty
    a_n x^n : a_0,a_1,\dots \in k\right\}$$
    (the sum does not need to converge). 
    It's not hard to show that $f(x) = a_0 + a_1x + \dots$
    is a unit if and only if $a_0\in k^\times$.
    Using this, we can show that all ideals of $k[[x]]$
    are in form $(0)$ and $(x^n)$ for $n\geq 0$.
    Only $(0)$ and $(x)$ are prime.
    \end{itemize}
\end{example}
In fact, a ring with unique maximal ideal
(like the last example)
is relatively simpler to study that it bears a special name.
\begin{definition}[Local ring]
    $A$ is a \vocab{local ring} if it has
    a unique maximal ideal.
\end{definition}
We also have that $k[[x_1,\dots,x_n]]$ is also local.
We will see later in \Cref{chap:localization} that 
the study of rings can be reduced to the study of local rings.
\begin{exercise}
    Prove that the bijection in the correspondence theorem
    (\Cref{thm:correspondence}) takes prime ideal to prime ideal
    and maximal ideal to maximal ideal.
\end{exercise}
%
%
%
%
%
\subsection{Zorn's Lemma and Krull's Theorem}
\begin{theorem}[Krull's Theorem]
\label{thm:krull}
\index{Krull's theorem}
Let $A$ be a ring. Then, any ideal $I\subsetneq A$
is contained in a maximal ideal. In particular,
maximal ideals exist.
\end{theorem}
This theorem is seemingly obvious because if $I=I_0$ 
is not maximal, then there exists a larger ideal 
$I_0\subsetneq I_1$.
This continues, and we get a chain of ideals
$$I_0\subsetneq I_1 \subsetneq I_2 \subsetneq\dots.$$
If ring $A$ is horrible (example: $\ZZ[x_1,x_2,\dots]$
and the chain is $(x_1)\subsetneq (x_1,x_2)\subsetneq \dots$), 
then this process might continue indefinitely.

The idea is that we need to ``go beyond infinity''.
Zorn's lemma is a tool from set theory that allows us to do that.

A \vocab{poset} (or \vocab{partially ordered set}) 
$(\calP,\leq)$ is a set $\calP$ with a partial ordering $\leq$.
The ordering should be reflexive ($a\leq a$),
antisymmetric ($a\leq b$ and $b\leq a$ implies $a=b$),
and transitive ($a\leq b$ and $b\leq c$ implies 
$a\leq c$).
A \vocab{totally ordered set} is a poset
whose any two elements are comparable.

The following lemma is equivalent to axiom of choice.
\begin{lemma}[Zorn's lemma]
    \label{lem:zorn}
    \index{Zorn's lemma}
    Let $(\calP,\leq)$ be a non-empty poset such that
    for any totally-ordered subset $\calT\subset \calP$,
    there exists some element $a_\calP\in \calT$
    such that $a_\calT\geq t$ for all $t\in \calT$.
    Then, $\calP$ has a maximal element,
    i.e, there exists $m\in \calP$ such that $m\geq p$
    for all $p\in \calP$.
\end{lemma}
\begin{proof}[Proof of \Cref{thm:krull}]
    Take a poset $(\calP, \subseteq)$, where
    $$\calP = \{\text{ideal } J : I\subseteq J\subsetneq A\}$$
    ordered by containment.
    
    We need to verify the condition for Zorn's lemma
    \Cref{lem:zorn}.
    Suppose that $\calT$
    is a totally-ordered subset of ideals,
    then we claim that $K = \bigcup_{J\in \calT} J$ is an ideal
    not equal to $A$. To do this, we need to check three things.
    \begin{itemize}
    \item \textbf{Closed under addition.}
    If $a,b\in K$, then there exists $J, J'\in\calT$
    such that $a\in J$ and $b\in J'$.
    Then, $a+b\in \max(J,J')$, which is a subset of $K$.
    \item \textbf{Absorbs multiplication.}
    If $a\in K$ and $r\in A$,
    then $a\in J$ for some $J\in \calT$,
    and so $ra\in J$, which implies $ra\in K$.
    \item \textbf{$K$ is not equal to $A$.}
    Note that $1\notin J$ for all $J\in \calT$,
    so $1\notin K$. 
    \end{itemize}
    Hence, $K$ is an element in $\calP$.
    It clearly contains every ideal in $\calT$,
    so the condition of Zorn's lemma is satisfied.
    Hence, poset $\calP$ has a maximal element,
    which must be a maximal ideal.
\end{proof}%
%
%
%
%
%
\section{Radicals}
\subsection{Nilradical}
Recall that $\Nil(A)$ is the set of all nilpotent
elements of $A$.
This is also called the \vocab{nilradical} of $A$.
We now use Zorn's lemma to prove the following 
important property of nilradical.
\begin{proposition}[Nilradical is intersection of prime ideals]
    \label{prop:nilradical} We have
    $$\Nil(A) = \bigcap_{\mathfrak p\subseteq A
    \text{ prime}}\mathfrak p$$ 
\end{proposition}
\begin{proof}
    We first show the direction $\subseteq$.
    To do this, we have to prove that
    $\Nil A\subset\mathfrak p$ for all prime ideal $\pp$.
    Indeed, $a\in\Nil A$, then $a^n=0\in\mathfrak p$,
    so $a\in\mathfrak p$ by definition of prime ideal.

    In the direction $\supseteq$,
    given $f\notin\Nil(A)$, we need to find a prime ideal
    not containing $f$.
    We use Zorn's lemma. Consider the poset $(\calP,\subseteq)$ where
    $$\calP = 
    \Big\{\text{ideal } I\subsetneq A\ :\ 1,f,f^2,\dots 
    \text{ are all not in } I
    \Big\}$$
    ordered by inclusion.
    
    Notice that $f\notin\Nil(A)$, so $S\neq\emptyset$,
    so $0\in S$. Moreover, for any totally ordered subset 
    $\calT\subseteq \calP$, it is not hard to check that
    $\bigcup_{J\in\calT} J$ is an ideal in $\calP$
    that dominates every element $\calT$.
    Thus, the condition of Zorn's lemma is satisfied,
    so $\calP$ has a maximal element $J$.
    Notice that $f\notin J$.

    Now, we claim that $J$ is prime
    (which will finish the proof).
    Assume $xy\in J$ but $x,y\notin J$.
    Then, $J+(x)\supsetneq J$, so $J+(x)$ must contain
    some power of $f$.
    In particular, $f^n\in J+(x)$ for some $n$
    and $f^m\in J+(y)$ for some $m$.
    However, by multiplying, we get that
    $$f^{m+n}\in (J+(x))(J+(y)) =  J+(xy)
    = J,$$ 
    which is a contradiction.
\end{proof}
%
%
%
\subsection{Radical Ideal}
We have shown that $\Nil(A)$ is the same as
the intersection of all prime ideals of $A$.
We generalize this a little bit, restricting to
only prime ideals that contains $I$.
\begin{definition}[Radical ideal]
    Given a ring $A$ and an ideal $I$,
    we define the \vocab{radical ideal}
    \notation{$\sqrt I$}{radical ideal}
    $$\sqrt I = \{a\in A : a^n\in I
    \text{ for some positive integer } n\}.$$
    (Example: if $A=\ZZ$, then $\sqrt{(48)} = (6)$.)
    
    An ideal $I$ is called \vocab{radical}
    if $\sqrt I = I$.
\end{definition}
Let $\pi : A\to A/I$ be the quotient map.
Then, the condition that $a^n\in I$ 
for some positive integer $n$
is equivalent to $\pi(a)\in\Nil(A/I)$. 
Therefore, we have 
\begin{align*}\sqrt I 
&= \pi^{-1}(\Nil(A/I)) \\
&= \pi^{-1} \left(\bigcap_{\pp\subseteq A/I
\text{ prime}} \pp\right) \tag{\Cref{prop:nilradical}} \\
&= \bigcap_{I\subseteq \pp\subseteq A\text{ prime}}
\pp. \tag{Correspondence theorem \Cref{thm:correspondence}}
\end{align*}
\section{Spectrum and Zariski Topology}
\label{sec:spec}
Several pictures were adapted from \cite{napkin}.
\subsection{Varieties}
\label{subsec:varieties}
Many aspects of commutative algebra are motivated 
by algebraic geometry. 
Since this is not algebraic geometry notes,
we will informally introduce algebraic geometry
just enough to understand some geometric intuition
behind commutative algebra, but we will not do 
any serious algebraic geometry.
\textbf{This entire subsection is very informal}
and is intended to only give intuition.
There is no actual math done in this subsection.

For these notes, a \vocab{variety} is a set $V$
of the form
$$V = \{\vec x\in\CC^n : 
f_1(\vec x) = f_2(\vec x) = \dots = f_m(\vec x)=0\}
\ \subseteq\ \CC^n,$$
where $f_1,\dots,f_m \in \CC[x_1,\dots,x_n]$.
To understand variety, we want to consider a \vocab{ring 
of regular functions} associated to it. 
More specifically, we want to consider a ring $A$ of 
all \textbf{regular functions} $f : V\to \CC$,
which are rational functions from $V\to\CC$
defined everywhere in $V$.
\begin{example}
    \begin{itemize}
    \item If $V=\CC^n$, then $A=\CC[x_1,\dots,x_n]$
    is our ring of regular functions. 
    Function $x_i$ would take a point in $V$ and output
    the $i$-th coordinate of $V$.
    Notice that even though we say
    rational function, we cannot have denominator here.
    \item Now, suppose that 
    $$V=\{(x,y)\in\CC^2 : x^2+y^2=1\}.$$
    Then, now, the function $x^2+y^2-1$ is the same as function $0$.
    Thus, the ring of regular functions 
    needs to be modded out by $(x^2+y^2-1)$,
    so $A = \CC[x,y] / (x^2+y^2-1)$.
    \end{itemize}
\end{example}
Thus, in general, if 
$V = \{\vec x\in\CC^n : 
f_1(\vec x) = f_2(\vec x) = \dots = f_m(\vec x)=0\}$,
then the ring of regular functions $A$ would be
$$\CC[x_1,\dots,x_n]/(f_1,\dots,f_m).$$

Many properties of ring of regular functions $A$ 
translates to properties
of varieties $V$. We describe some connections below,
although none of these are precise.
\begin{itemize}
\item Given an ideal $I\subset A$, the set of points
in $V$ whose every function in $I$ vanishes
is a closed subset of $V$ (in analytical sense).
\item Suppose there is another variety $V'$
with ring of regular functions $A'$.
Then, a morphism $\phi : V\to V'$
corresponds to a ring morphism $\phi^* : A'\to A$
that takes a function $f : V'\to \CC$ 
to a function $f\circ\phi : V\to \CC$.
\item $A$ is a domain if and only if $V$ is irreducible
(i.e., is not a union of two subvarieties).
Suppose that $V$ is a union of two subvarieties $V_1\cup V_2$,
where $V_1$ is cut out by equations $f_1(\vec x)
= \dots = f_m(\vec x)=0$, while
$V_2$ is cut out by equations
$g_1(\vec x)= \dots = g_n(\vec x)=0$.
Then, $V_1\cup V_2$ is cut out by equations 
$f_i(\vec x)g_j(\vec x) = 0$.
In particular, the function $f_i$ and $g_j$
are nonzero regular functions on $V$,
but they multiply to $0$, which makes $A$ not a domain.
\item $A$ is a field if and only if $V$ has exactly one point.
This is because informally, if $V$ has more than one point,
we can create a function $f$ vanishing at one point
but not the other, and this $f$ would not be invertible
and not a unit.
\item For the definition of variety defined above,
we always have that $A$ is reduced.
In particular, cutting out by equation $(x-y)^2=0$
or $x-y=0$ mean the same thing. 
We can have a more subtle notion of geometric objects
where $\CC[x,y]/((x-y))$ is simply a line $x=y$,
but $\CC[x,y]/((x-y)^2)$ carries more information
than a line, including an infinitesimal neighborhood
around each point in the line $x=y$.
This infinitesimal neighborhood is informally called a fuzz.

This notion of geometric object that capture fuzz
is a \vocab{scheme}, which generalize varieties in 
a few different directions, one of which is allowing fuzz.
We do not study schemes in these notes.
\end{itemize}

One idea of algebraic geometry is to go backward:
given a ring of function $A$,
construct a geometric space that can substitute
the variety $V$ above. The space is called $\Spec A$,
and it will make all connections above precise.
\subsection{Spectrum of a Ring}
\begin{definition}[Spectrum]
    \begin{itemize}
    \notation{$\Spec A$}{spectrum of $A$}
    \item The \vocab{spectrum} of a ring $A$,
    denoted $\Spec A$ is the set of
    all prime ideals of $A$.
    \notation{$\mSpec A$}{maximal spectrum of $A$}
    \item The \vocab{maximal spectrum} of a ring $A$
    is denoted $\mSpec A$
    and is the set of all maximal ideals of $A$.
    \end{itemize}
\end{definition}
For example, by Hilbert Nullstellensatz
(\Cref{thm:max_ideal_poly_ring}),
maximal ideal in $\mSpec\mathbb C[x_1,\dots,x_n]$
correspond to points in $\mathbb C^n$.
Maximal ideal $\mm = (x_1-a_1,\dots,x_n-a_n)$
corresponding to point $(a_1,\dots,a_n)\in\CC^n$,
and for any polynomial $f$, we have that 
$f\bmod \mm = f(a_1,\dots,a_n)$ is the value of 
$f$ evaluated at that point.

In general, we can view a ring as a ring of regular functions.
For any ring $A$ and $f\in A$,
we want to think of $f$ as defining a ``function''
on $\mSpec A$.
This generalizes to $\Spec A$, and we will explain 
later (in \Cref{rmk:why_prime}) 
why we prefer to use prime ideals instead of maximal ideal.
\begin{definition}
    Given $\mathfrak p\in \Spec A$ and $f\in A$,
    the ``value'' of $f$ when ``evaluated'' at $\mathfrak p$
    is $f\bmod\mathfrak p \in A/\mathfrak p$.
\end{definition}
For psychological reason, we occasionally use a different notation for
points in $\Spec A$ with the prime ideals of $A$. Specifically,
\begin{itemize}
    \item Point $x\in \Spec A$ corresponds to prime ideal 
    $\pp_x\subseteq A$.
    \item Prime ideal $\pp\subseteq A$ corresponds to point $[\pp]\in\Spec A$.
\end{itemize}

\subsection{Zariski Topology}
We now equip a topology of $\Spec A$,
which will be a substitute topology for our geometric space.

As a quick reminder, to define a topological space $X$,
one needs the set of \vocab[open sets]{open} sets $\mathcal U$
preserved by unions and finite intersections.
A set is \vocab[closed sets]{closed} iff its complement is open.
Such collection is very arbitrary,
and many intuition from metric spaces will break down.
\begin{definition}[Zariski topology]
    Let $A$ be a ring. For each ideal $I\subseteq A$, define 
    \notation{$V(I)$}{closed set of Zariski topology}
    $$V(I) = \{\pp\in\Spec A : I\subseteq \pp\}.$$
    In \vocab{Zariski topology} of $\Spec A$,
    a subset $S\subseteq\Spec A$
    is closed if and only if $S=V(I)$ for some ideal $I\subseteq A$.
\end{definition}
We have to check that this defines a topology.
We verify the following.
\begin{itemize}
    \item $\emptyset = V((1))$, so $\emptyset$ is closed.
    \item $\Spec A = V((0))$, so $\Spec A$ is closed.
    \item $\bigcap_\alpha V(I_\alpha) 
	= V\left(\sum_\alpha I_\alpha\right)$,
	so the set of $V(I)$'s is indeed closed under intersection.
    \item $V(I)\cup V(J) = V(IJ)$,
	so the set of $V(I)$'s is closed under finite unions.
\end{itemize}
The last one is the least trivial.
In one direction, $IJ\subset I$, so $V(IJ)\supset V(I)$,
and so $V(IJ)\supset V(I)\cup V(J)$.
Now, if $x\in V(IJ)$ but $x\notin V(I), V(J)$,
then there exists $i\in I, j\in J$ such that
$i,j\notin\mathfrak p_x$.
However, this implies that $ij\in IJ$ but
$ij\notin\mathfrak p_x$. This forms a contradiction.

The idea of this topology is that we want the coarsest 
(having fewest open sets) topology
such that functions in $A$ is continuous.
This means we need preimage of $\{0\}$ 
with respect to $f$ to be closed. This preimage is $V(f)$,
which we then declare that it's closed.
Taking intersection across many $f$'s gives that
we must declare $V(I)$ to be closed for all ideal $I$.
This is all the closed sets we need.

\begin{exercise}
    Prove that $\mm$ is a maximal ideal if and only if
    the set $\{[\mm]\}$ consisting of a single point
    $[\mm]$ is closed.
    Thus, maximal ideals correspond to \vocab{closed points}.
\end{exercise}

Now, let's get into some examples of drawing spectrum.
\begin{example}[$\Spec k$]
    If $k$ is a field, $\Spec k$ contains only one point.
\end{example}
\begin{example}[$\Spec \CC[x]$]
    Since prime ideals of $\CC[x]$ are $(0)$
    and $(x-\alpha)$ for each $\alpha\in \CC$,
    points in $\Spec \CC[x]$ are
    \begin{itemize}
	\item the \vocab{generic point} $(0)$ and
	\item points $(x-\alpha)$ for each $\alpha\in\CC$.
    \end{itemize}
    The closed sets of $\Spec\CC[x]$ are finite subsets
    of $\{(x-\alpha) : \alpha\in\CC\}$.
    Any open set must contain $(0)$ and all but 
    finitely many points in $\{(x-\alpha) : \alpha\in\CC\}$.
    \begin{center}
    \begin{asy}
    size(8cm);
    pair A = (-8,0); pair B = (8,0);
    draw(A--B, blue, Arrows);
    dot("$(x)$", (0,0), dir(-90), blue);
    dot("$(x-3)$", (3,0), dir(-90), blue);
    dot("$(x+i)$", (-4,0), dir(-90), blue);
    dot("$(0)$", (0,1), dir(90), red);
    draw( (-6,0)..(-3,0.5)..(-1,1)--(1,1)--(3,0.5)--(6,0), dotted+red );
    label("$\operatorname{Spec} \mathbb C[x]$", A, dir(-90), blue);
    \end{asy}
    \end{center}
    Point $[(0)]$ is more interesting.
    The set $\{[(0)]\}$ is not closed,
    but its closure is the entire $\Spec A$.
\end{example}
\begin{example}[$\Spec\ZZ$]
    $\Spec\ZZ$ looks exactly the same as the above picture.
    The points are
    \begin{itemize}
    \item the generic point $(0)$.
    \item the points $(p)$ for prime $p$.
    \end{itemize}
    \begin{center}
    \begin{asy}
    size(8cm);
    pair A = (-8,0); pair B = (8,0);
    draw(A--B, blue, Arrows);
    dot("$(7)$", (0,0), dir(-90), blue);
    dot("$(3)$", (3,0), dir(-90), blue);
    dot("$(19)$", (-4,0), dir(-90), blue);
    dot("$(0)$", (0,1), dir(90), red);
    draw( (-6,0)..(-3,0.5)..(-1,1)--(1,1)--(3,0.5)--(6,0), dotted+red );
    label("$\operatorname{Spec} \mathbb Z$", A, dir(-90), blue);
    \end{asy}
    \end{center}
\end{example}

Another way to describe the topology of $\Spec A$
is by giving open set.
\begin{proposition}[Distinguished open sets]
    \notation{$D(f)$}{distinguished open sets of 
    Zariski topology}
    For each $f\in A$, define the \vocab{distinguished open set}
    $$D(f) = \{\pp : f\notin \pp\} \ \subseteq\ \Spec A.$$
    Then, $\{D(f) : f\in A\}$ forms a basis of topology
    for $\Spec A$. In other words, any open set of $\Spec A$
    is a union of set of the form $D(f)$.
\end{proposition}
\begin{proof}
    Exercise.
\end{proof}
\subsection{Contraction and \texorpdfstring{$\Spec$}{Spec} Functor}
We now have a map that takes a ring $A$ to a topological space $\Spec A$.
Now, we want to promote this to a functor.
In order to do so, we need the following proposition.
\begin{proposition}[Contraction of prime ideal is prime]
    If $A,B$ be rings and $\phi : A\to B$ be a homomorphism.
    If $\mathfrak q\subseteq B$ is a prime ideal,
    then the contraction $\phi^{-1}(\mathfrak q)$
    is prime.
\end{proposition}
\begin{proof}
    Suppose $a,b\in A$ such that $ab\in \phi^{-1}(\frakq)$.
    Then, $\phi(ab)\in\qq$, so $\phi(a)\phi(b)\in\qq$.
    Since $\qq$ is prime, we have that 
    $\phi(a)\in\qq$ or $\phi(b)\in\qq$.
    Thus, $a\in\phi^{-1}(\qq)$ or $b\in\phi^{-1}(\qq)$,
    proving that $\phi^{-1}(\qq)$ is prime.
\end{proof}
\begin{remark}
    \label{rmk:why_prime}
    On the other hand, contraction of maximal ideals
    need not be maximal ideals.
    For example, if $A=\ZZ$, $B=\QQ$, $\phi$
    is the inclusion map, and $\qq=(0)$,
    then $\phi^{-1}\qq = (0)$, which is not a maximal ideal.
    This is the reason why we prefer to include all 
    prime ideals instead of using only maximal ideals.
    It is so that we can promote every ring morphism
    $\phi : A\to B$ to morphism of $\phi^\# : \Spec B\to \Spec A$.
\end{remark}
The above proposition defines a map from
$\Spec B\to \Spec A$ by 
$[\mathfrak q]\mapsto [\phi^{-1}(\mathfrak q)]$.
This induces a \textbf{contravariant functor}.
In particular, given a ring morphism $\phi : A\to B$,
we get a morphism $\phi^\# : \Spec B\to\Spec A$
of spectrum in the opposite direction.
\notation{$\phi^\#$}{induced map on $\Spec$}
\begin{exercise}
    Show that $\phi^\#$ is continuous (with respect to 
    Zariski topology).
\end{exercise}
\begin{example}
The quotient projection map $\pi : A\to A/I$.
induces the map $\pi^\# : \Spec A/I \to \Spec A$.
Map $\pi^\#$ is a homeomorphism from $\Spec A/I$
to a closed subset $V(I)$ of $\Spec A$.
\end{example}
\begin{exercise}
    The inclusion map $\RR[x]\to\CC[x]$ 
    induces a map $\Spec\CC[x]\to \Spec\RR[x]$.
    Draw picture of $\Spec\RR[x]$ and $\Spec\CC[x]$
    that visualizes this map.
\end{exercise}
Now, go back to \Cref{subsec:varieties} 
and convince yourself that the space $\Spec A$
substitute variety $V$ and still makes the connection
outlined in the bulleted list valid.
\subsection{\texorpdfstring{$\Spec\mathbb Z[x]$}{Spec Z[x]} and 
\texorpdfstring{$\Spec\mathbb C[x,y]$}{Spec C[x,y]}}
We draw pictures of $\Spec \ZZ[x]$ and $\Spec\CC[x,y]$.
\begin{theorem}[Ideals of $A[y]$ where $A$ is a PID]
    \label{thm:ideal_pid}
    Let $A$ be a PID. Then, the prime ideals of $A[y]$ are
    \begin{itemize}
    \item the zero ideal $(0)$;
    \item the ideal $(f(y))$, where $f$ irreducible;
    \item pick an irreducible element $p\in A$ and $g(y)\in A[y]$
	such that $\overline g(y) \in A/(p)[y]$
	is irreducible (in field $A/(p)$),
	then we have $(p, g(y))$.
    \end{itemize}
    Only the last category is a maximal ideal.
\end{theorem}
\begin{proof}
We leave this as \Cref{prob:ideal_pid}.
\end{proof}
\paragraph{$\Spec \CC[x,y]$}
Now let's try to draw $\Spec \CC[x,y]$ (i.e., $A=\mathbb C[x]$).
We have the ideals $(0)$ and $(f(x,y))$
for the first two categories.
For the last one, all prime elements must be in form
$p=(x-\alpha)$, and $A/(p) \simeq \mathbb C$.
Thus, an irreducible polynomial in
$A/(p)[y] = \mathbb C[y]$ must be in form $y-\beta$,
giving ideals $(x-\alpha, y-\beta)$
for $(\alpha, \beta)\in\mathbb C^2$.
Therefore, $\Spec\CC[x,y]$ looks like the following:
\begin{itemize}
    \item The ideal $(x-\alpha, y-\beta)$,
	which form a plane $\mathbb C^2$.
    \item The ideal $(f)$ for each irreducible $f\in\CC[x,y]$,
    which will look like a cloud above 
    the irreducible curve corresponding to .
    \item The generic point, which look like a cloud above everything.
\end{itemize}
Here is a picture that shows three types of points.
\begin{center}
\begin{asy}
size(6cm);
import graph;
real f(real x) { return x*x; }
graph.xaxis("");
graph.yaxis("");
draw(graph(f,-2,2,operator ..), red+dashed, Arrows(TeXHead));
dot("$(y-x^2)$", (1.3, f(1.3)), dir(-45), red);
dot("$(x-1,y+2)$", (1,-2), dir(-45), blue);
pair O = (-3,3);
dot("$(0)$", O, dir(225), red);
filldraw(circle(O, 0.8), opacity(0.2)+lightred, dashed+red);
\end{asy}
\end{center}
\paragraph{$\Spec \ZZ[x]$}
$\Spec\mathbb Z[x]$ (i.e., $A=\ZZ$) looks pretty similar,
except that the horizontal axis will be
replaced by primes.
The points are as follows.
\begin{itemize}
    \item Maximal ideals, which are of the form $(p,f(x))$
    where $f\in\FF_p[x]$.
    \item For each irreducible polynomial $f\in\ZZ[x]$,
    the point $(f)$, which is a cloud passing through
    $(p,\overline f(x))$ for each prime $p$.
    \item For each prime $p$, the ideal $(p)$,
    which is a cloud above $(p,f(x))$
    for any polynomial $f$.
    \item The generic point $(0)$, which is a cloud
    above everything.
\end{itemize}
\begin{center}
\begin{asy}
size(12cm);
defaultpen(fontsize(9pt));
int[] p = {2,3,5,7};

void cloud(pair P, real r, string s, pair l){
	filldraw(circle(P, r), red+opacity(0.1), red+dashed);
	dot(P, red);
	label(s, P + r*l/abs(l), l);
}
cloud((3.7,2.4), 0.3, "$(0)$", NE);

draw((-0.3,0) -- (4.5,0), blue);
cloud((4.2,0), 0.2, "$(x)$", NE);

draw((-0.3,0.9) .. (0,1) .. (1,4/3) .. (2,8/5) .. (3, 12/7) 
.. (4.2,2) .. (4.5,2), blue);
cloud((4.2,2), 0.2, "$(x-1)$", NE);

draw((4.5,1.2) .. (4.2,1.2)..(3,1.01){W} .. (2,1.2) .. (1,1.01){W} .. (0.5,1.2) ..
(0,1) .. (0.5,0.8) .. (1,0.99){E} .. (2,0.8) .. (3,0.99){E}
.. (4.2,0.8), blue);
cloud((4.2,1.2), 0.2, "$(x^2+1)$", NE);

pair[][] label_dir = {{NE, E}, {NE, E, E}, {NE, E, SE, NE, NE},
{NE, E, E, E, E, E, SE}};
for(int i=0; i<4; ++i){
	draw((i,-0.2) -- (i,2.7));
	for(int j=0; j<p[i]; ++j){
		if((i != 3) || (j == 0) || (j==6))
			label("$(" + ((string)p[i]) + ", x+" + ((string) j) + ")$", (i, 2*j/p[i]), 
			label_dir[i][j]);
		dot((i,2*j/p[i]), red);
	}
	cloud((i,2.3), 0.2, "$(" + ((string)p[i]) + ")$", NE);
	label("$V(" + ((string)p[i]) + ")$", (i,-0.2), S);
}
dot((1,1), red);
label("$(3,x^2+1)$", (1,1), E);
dot((3,1), red);
label("$(7,x^2+1)$", (3,1), E);
\end{asy}
\end{center}
\section{Problems}
\begin{problem}
Let $A$ be a ring.
\begin{alphenum}
\item Is it true that $I(J+K) = IJ + IK$ for all ideals $I,J,K$ of $A$?
\item Is it true that $I(J\cap K) = IJ\cap IK$ for all ideals $I,J,K$ of $A$?
\end{alphenum}
\end{problem}
\begin{problem}
Let $A$ be a ring and ideals $I$ and $J$ are coprime
(i.e., $I+J=A$). Prove that for all positive integers $m,n$,
ideals $I^m$ and $J^n$ are coprime.
\end{problem}
\begin{problem}
Let $A$ be a ring such that for all $x\in A$, there exists an integer $n>1$ (possibly depending on $x$) such that $x^n=x$. Prove that every prime ideal of $A$ is maximal.
\end{problem}
\begin{problem}
Let $A$ be a ring.
Using Zorn's lemma, prove that $A$ has a 
\vocab{minimal prime ideal},
a prime ideal minimal under inclusion.
\end{problem}
\begin{problem}[Prime Avoidance, \textbf{to be used later}]
\label{prob:prime_avoidance}
\index{prime avoidance}
Let $A$ be a ring and $\mathfrak p_1,\dots,\mathfrak p_n$ be prime ideals. Suppose that $I$ is an ideal contained in $\bigcup_{i=1}^n\mathfrak p_i$. Prove that $I\subseteq\mathfrak p_i$ for some $i$.
\end{problem}
\begin{problem}
\label{prob:ideal_pid}
Prove \Cref{thm:ideal_pid}.
\end{problem}
\begin{problem}
Let $A$ be a ring. Prove that the following are equivalent.
\begin{enumerate}[label=(\alph*)]
\item $A$ has exactly one prime ideal.
\item Every element of $A$ is either nilpotent or a unit.
\item $A/\operatorname{Nil} A$ is a field.
\end{enumerate}
\end{problem}
\begin{problem}
Prove that $\Spec A$ is quasicompact (\Cref{def:compact}) for any ring $A$.
\end{problem}
\begin{problem}
\label{prob:product_ring}
Let $A$ be a ring. Prove that the following are equivalent.
\begin{enumerate}[label=(\alph*)]
\item $\operatorname{Spec} A$ is a disjoint union of two closed sets $X_1$ and $X_2$.
\item There exists nonzero elements $e_1, e_2\in A$ such that 
$e_1+e_2=1$ and $e_1e_2=0$.
\item $A = A_1\times A_2$ for some rings $A_1$ and $A_2$.
\end{enumerate}
Moreover, when $A=A_1\times A_2$,
then $\Spec A$ is homeomorphic to disjoint union 
$\Spec A_1$ and $\Spec A_2$.
\end{problem}
\begin{problem}
Is $\Spec \prod_{n=1}^\infty \FF_2$ homeomorphic 
to the disjoint union $\coprod_{n=1}^\infty \Spec\FF_2$?
\end{problem}